\section{Data management}\label{sec:data_management_python}
\subsection{Data format handled}\label{sec:data_format_handled}
\gretlpy expects a dataset provided in one of the following formats:
\begin{enumerate}
	\item \texttt{gdt} or \texttt{gdtb} (\lgretl['s] native formats),
		\texttt{csv}, \texttt{dta} (\pkg{Stata}), \texttt{wf1} (\pkg{Eviews}),
		or spreadsheets (\texttt{xls}, \texttt{xlsx}, and \texttt{ods}).
	\item as a \python['s] dictionary with the following entries
	(see \ref{sec:gretl4py:data:dict} for details):
	\begin{lstlisting}[style=pythonstylefoot]
{
"varname":     <variable names (list)>,
"obs_labels":  <observation labels (list)>,
"data_array":  <observations (NumPy array)>,
"periodicity": <dataset periodicity (int)>
}
	\end{lstlisting}
	Such dictionary can be created manually or with the help
	of one of the utility functions described in
	Section~\ref{sec:data_management_python}.
\end{enumerate}
Note: If the provided dataset contains a variable named
`const', it will automatically be renamed to `const\_'.


\subsection{Loading a dataset}
To load datasets, \gretlpy provides the \pythcodeidx{get\_data()}
function (see \ref{get_data()}). Using this function, data in all
supported formats may be loaded, and it is the only supported
interface for dataset loading (that is, direct calls to the
\classpydataset constructors are not allowed). In the simplest
case, i.e.\ when a user wants to load a dataset bundled%
\footnote{All datasets bundled with each \gretlpy distribution
	are listed~\error{XXX}.}
with a \gretlpy build, it is sufficient to run the following snippet%
\footnote{Note that the call \pythcode{gretl.get_data('bjg.gdt')}
	is equivalent to
	\pythcode{gretl.get_data(gretl.gretldir + '/data/bjg.gdt')}.
	This is because \texttt{gretl.gretldir + '/data/'} is included
	in the \lgretl search path.}:
\begin{lstlisting}[style=pythonstylefoot]
import gretl

d = gretl.get_data('bjg.gdt')
\end{lstlisting}

Here, \texttt{bjg.gdt} denotes the monthly dataset on international
airline passengers used by Box and Jenkins in \textit{Time Series
	Analysis: Forecasting and Control}
to illustrate multiplicative seasonal MA models (see
\url{http://www.itl.nist.gov/div898/handbook/pmc/section4/pmc44a.htm}
for details).


\subsection{Default dataset}\label{sec:globalandindivodual()}
The \gretlpy package may handle an arbitrary number of datasets,
one of which may be set as the default. The effect is that whenever
a user calls a function or method that allows passing a dataset
via the \pythcode{data} argument, and no dataset is provided in
that way, the procedure in question automatically uses the
currently defined default dataset.

The default dataset is defined according to the following rules:
\begin{itemize}
	\item If a dataset is successfully loaded and no default
		dataset has been defined yet, the current dataset is set
		as the default. As a consequence, the dataset loaded first
		is always set as the default.
	\item The user explicitly calls the \pythcodeidx{set\_as\_default()}
		method (see \ref{met:gretldataset:set_as_default}).
\end{itemize}

The simplest way to verify whether a given dataset is set as the
default is to check the value of the \pythcodeidx{is\_default}
attribute (see \ref{attr:gretldataset:is_default}). Alternatively,
the user may call the \pythcodeidx{show\_data\_status()} function,
which prints all available \classpydataset instances (see
\ref{show_data_status}).

For details on the functions and methods for data management
provided by \gretlpy, please refer to Section~\ref{sec:classpydataset}.
